\chapter{Compact Hall Geometry on a K3 Product}
\label{cycompact:chapter}

Translation along an elliptic curve removes a common position while
preserving the positions of objects relative to one another. An extension
depends on these relative positions: sheaves in different fibres have
no extensions, whereas sheaves in the same fibre can have both
deformations and obstructions. A Hall construction on the translation
quotient must therefore retain one elliptic torsor throughout its flag
correspondences.

We develop this geometry for $X=S\times E$, where $S$ is a smooth
projective complex K3 surface and $E$ is an elliptic curve. The disjoint
point locus has explicit self-dual coefficients and finite trace maps.
Compact rational-curve pairs give an oriented critical model with
computable cohomology. At a polystable collection, cyclic deformation
theory gives one formal model for all finite flags. These constructions
also expose the data needed at collisions: scalar trace complexes,
paired orbit directions, square-root sewing and supported coefficient
maps. Their compatibility on the full compact extension stack is the
global comparison problem.

Fix nonzero holomorphic forms $\omega_S\in H^0(S,K_S)$ and
$\omega_E\in H^0(E,K_E)$, and use their product as the volume form
on $X$. Coherent complexes are over $\mathbb C$; constructible
complexes have rational coefficients. All shifts are cohomological,
so $V[a]^j=V^{j+a}$. Derived pullbacks and tensor products retain
this convention. Put $W=T_0E$.

\section{One translation torsor for a whole flag}
\label{cycompact:sec:torsor}

Let $\mathcal M_\alpha$ be an $E$-invariant moduli stack of sheaves
of charge $\alpha$, and let $\mathcal Q_\alpha=[\mathcal M_\alpha/E]$.
The morphism $\mathcal Q_\alpha\to BE$ remembers the principal
$E$-bundle in the definition of the quotient stack. An object of
$\mathcal Q_\alpha$ over a base $B$ is thus an $E$-torsor
$P\to B$ together with an equivariant family on $P$.

\begin{proposition}\label{cycompact:prop:torsor}
There is a natural equivalence
\[
 \mathcal Q_\alpha\times_{BE}\mathcal Q_\beta
 \simeq[(\mathcal M_\alpha\times\mathcal M_\beta)/E_{\mathrm{diag}}].
\]
Exact sequences on the common torsor give a correspondence
\begin{equation}\label{cycompact:eq:span}
 \mathcal Q_\alpha\times_{BE}\mathcal Q_\beta
 \xleftarrow{\ a\ }\mathcal H_{\alpha,\beta}
 \xrightarrow{\ p\ }\mathcal Q_{\alpha+\beta}.
\end{equation}
Finite filtrations identify all successive compositions of these
correspondences, compatibly with refinement. The zero object gives
the relative unit $BE$.
\end{proposition}

\begin{proof}
An object of the fibre product consists of two torsors, their
identification, and the two equivariant families. Transport along
the identification makes this one torsor with a pair of families.
This is precisely the diagonal quotient, including morphisms and
descent over $B$. An exact sequence in that family defines
\eqref{cycompact:eq:span}.

For three successive quotients, either composite correspondence
classifies a filtration $0\subset F_1\subset F_2\subset F_3$ on
one torsor. A further refinement inserts further subobjects into
that same filtration. Forgetting specified levels is independent
of the order in which the other levels are forgotten. This proves
the assertions for every finite flag. The zero sheaf is fixed by
$E$, so its quotient is $BE$; inserting it makes no additional
choice.
\end{proof}

The relative product is already necessary for a single point.
Let $\mathfrak M_1$ denote the stack of flat length-one sheaves on
$X$. Then
\begin{equation}\label{cycompact:eq:point-stack}
 \mathfrak M_1=X\times B\mathbb G_m,\qquad
 [\mathfrak M_1/E]\times_{BE}[\mathfrak M_1/E]
 \simeq S^2\times E\times B\mathbb G_m^2.
\end{equation}
To see the first statement in families, the support of a flat
length-one sheaf is finite over the base. Its pushforward is a
line bundle $L$. The support algebra acts faithfully on $L$,
and its action identifies it with the scalar algebra of the base.
The support is therefore the graph of a map to $X$, and the sheaf
is that graph's pushforward of $L$. This construction respects
isomorphisms and descent. For the second statement, write the two
elliptic coordinates as $(e_1,e_2)$ and quotient by their common
translation. The surviving coordinate is $e_2-e_1$.

\section{Disjoint point coefficients}
\label{cycompact:sec:disjoint}

Let $U_n\subset X^n$ be the variety of ordered pairwise distinct
points. The torus $T_n=(\mathbb G_m)^n$ acts trivially and records
the automorphisms of the coefficient lines. The symmetric group
$\mathfrak S_n$ permutes points and torus factors. Put
\[
 \mathfrak Q_n=[U_n/(E\times(T_n\rtimes\mathfrak S_n))],
 \qquad \mathfrak Q_0=BE.
\]
After trivializing the elliptic torsor, the morphism
$q_n:\mathfrak Q_n\to BE$ becomes the smooth stack
$[U_n/(T_n\rtimes\mathfrak S_n)]$, of complex dimension
$3n-n=2n$. Complex orientation consequently gives
\begin{equation}\label{cycompact:eq:disjoint-coefficients}
 \omega_{q_n}=\mathbb Q_{\mathfrak Q_n}[4n],\qquad
 C_n=\mathbb Q_{\mathfrak Q_n}[2n],\qquad
 C_n\otimes C_n\simeq\omega_{q_n}.
\end{equation}
Here $\omega_{q_n}$ is the relative dualizing complex. Thus $C_n$
is self-dual for relative Verdier duality and $C_0=\mathbb Q_{BE}$.
These formulas specify smooth coefficients on the disjoint locus.

Let $D_{m,n}\subset\mathfrak Q_m\times_{BE}\mathfrak Q_n$ be
the open locus of disjoint supports, and restrict the extension
correspondence to
\[
 D_{m,n}\xleftarrow{\ a\ }\mathfrak F_{m,n}
 \xrightarrow{\ p\ }\mathfrak Q_{m+n}.
\]

\begin{proposition}\label{cycompact:prop:disjoint}
The map $a$ is an equivalence, and $p$ is representable finite
\'{e}tale of degree $\binom{m+n}{m}$. There is a canonical
identification
\[
 a^*(C_m\boxtimes_{BE}C_n)=p^*C_{m+n}.
\]
The adjunction and finite trace maps have composite
\[
 C_{m+n}\longrightarrow p_*p^*C_{m+n}
 \xrightarrow{\operatorname{Tr}_p}C_{m+n},
 \qquad \binom{m+n}{m}\operatorname{id}.
\]
They are adjoint for the relative pairings, and their refinement
identities hold for every finite flag, including empty blocks.
\end{proposition}

\begin{proof}
Derived Hom vanishes between sheaves with disjoint support.
Every extension therefore splits uniquely. A subobject of length
$m$ in a fixed middle sheaf chooses an $m$-element subset of its
support. After an \'{e}tale ordering of that support, $p$ is the
constant cover by those subsets. This proves both geometric
assertions. The coefficient shifts agree because
$2m+2n=2(m+n)$.

On that cover, the first map repeats a section on every sheet,
and trace sums its values. Their composite and adjunction follow
from summing $\binom{m+n}{m}$ identical values and from
$\sum_I uv_I=u\sum_Iv_I$. A refined flag is an ordered partition
of the support. Every parenthesization sums over the same
partitions, which proves all refinements at once. An empty block
contributes one choice and the coefficient $C_0$.
\end{proof}

The scalar automorphisms also enter the determinant line. After
fixing a basis of $W$, its pullback to $U_n$ is
$L_n=K_{U_n}\otimes\det\mathbb C^n$. A permutation acts by
$\operatorname{sgn}^3$ on the support volume and by
$\operatorname{sgn}$ on the torus determinant. Their product is
one. If $\lambda_n$ places the support factors before the scalar
factors, binary sewing has sign $(-1)^{mn}$. The generators
\begin{equation}\label{cycompact:eq:disjoint-sign}
 \widetilde\lambda_n=(-1)^{n(n-1)/2}\lambda_n
\end{equation}
sew with coefficient one, since
$\binom{m+n}{2}=\binom m2+\binom n2+mn$.
The same identity gives every associative regrouping.

\subsection{A two-sheet coefficient cannot be replaced by a scalar}

Choose a nonzero two-torsion point $\tau\in E$. Translation by
$\tau$ preserves the sheaf supported at $(s,e)$ and $(s,e+\tau)$,
and interchanges its two length-one subobjects. On that fibre
the flag coefficient is $A=\mathbb Q^2$, with coordinatewise
product and the transposition action. Its trace is
$\operatorname{Tr}(a,b)=a+b$.

\begin{proposition}\label{cycompact:prop:two-sheet}
There is no equivariant unital algebra map $A\to\mathbb Q$ when
the target has the trivial action. The complex
$[A\xrightarrow{\operatorname{Tr}}\mathbb Q]$ in degrees zero
and one has only one cohomology group: the sign representation
in degree zero.
\end{proposition}

\begin{proof}
A unital algebra map sends the two orthogonal idempotents
$(1,0),(0,1)$ to $0,1$ in some order. It therefore cannot be
invariant under their transposition. Trace is surjective and
has kernel $\mathbb Q(1,-1)$, on which the transposition acts
by $-1$.
\end{proof}

This kernel is not an ideal: $(1,-1)^2=(1,1)$ has trace two.
Nor does adding a homologically scalar dg algebra avoid the
multiplicity. Indeed, let $B$ be a unital dg algebra with
$H^0(B)=\mathbb C$, and let a finite group $G$ act transitively
on a nonempty finite set $\Omega$. Let $G$ act on $M_m(B)$ by
conjugation through a representation on $\mathbb C^m$, trivially
on $B$. A $G$-equivariant unital dg algebra map
$\mathbb C^\Omega\to M_m(B)$ induces orthogonal projectors in
$M_m(\mathbb C)$. Transitivity makes their ranks equal. Since
they sum to the identity, $m=d|\Omega|$ for an integer $d\geq1$.
The permutation representation and its diagonal projectors
attain $m=|\Omega|$.

\section{Trace complexes on split and nonsplit flags}
\label{cycompact:sec:trace}

The structure-sheaf cohomology of the product has the finite
Dolbeault model
\begin{equation}\label{cycompact:eq:scalar-model}
 H=\mathbb C[\sigma]/(\sigma^2)\otimes\Lambda(\eta),
 \qquad |\sigma|=2,\quad |\eta|=1,\quad
 \int_X\eta\sigma=1.
\end{equation}
Parallel representatives for a product Ricci-flat metric give
the multiplicative inclusion into the Dolbeault complex used in
Chapter~\ref{cyquant:chapter}. There is one cohomology class in
each degree zero through three, and multiplication followed by
the displayed integral is a perfect pairing of degree $-3$.

For a vector bundle $A$ write
$C_A=R\Gamma(X,R\mathcal Hom(A,A))$ and let $C_A^0$ be its
trace-free summand. Positive rank is needed for the scalar
splitting $h\mapsto h\operatorname{id}_A/\operatorname{rk}A$.
Write $C(A,B)=R\Gamma(X,R\mathcal Hom(A,B))$.

\begin{proposition}\label{cycompact:prop:trace-split}
For bundles of positive ranks $n,m$ there is a decomposition
\begin{equation}\label{cycompact:eq:trace-split}
 C_{A\oplus B}^0\simeq
 C_A^0\oplus C_B^0\oplus C(A,B)\oplus C(B,A)\oplus H.
\end{equation}
The last summand is
$h\mapsto\operatorname{diag}(h/n,-h/m)$, where the two entries
multiply the respective identities. Its induced pairing is
\[
 (h,k)\longmapsto\frac{n+m}{nm}\int_Xhk.
\]
The construction applies to relative endomorphism complexes in
families of bundles, with the relative structure-sheaf complex.
\end{proposition}

\begin{proof}
Split the endomorphism matrix into its four blocks. The two
diagonal traces must sum to zero, leaving one copy of the
scalar complex. Composition of their scalar matrices has
trace $(1/n+1/m)hk$. All decompositions and maps are defined
before taking cohomology, so they also give the relative claim.
\end{proof}

For $r$ bundles of ranks $n_i>0$, put
\[
 U_r=\{(a_1,\ldots,a_r)\in\mathbb C^r:
                 \textstyle\sum_i a_i=0\}.
\]
The scalar complex is $H\otimes U_r$. In its trace
coordinates $h_i\in H$, its pairing is
\[
 \langle h,k\rangle=\int_X\sum_i\frac{h_ik_i}{n_i}.
\]
The scalar matrix in the $i$th block is $h_i/n_i$. For a
partition into nonempty groups $I_1,\ldots,I_s$, set
$n_{I_a}=\sum_{i\in I_a}n_i$. The maps
\begin{equation}\label{cycompact:eq:trace-grouping}
 G(h)_a=\sum_{i\in I_a}h_i,\qquad
 J(k)_i=\frac{n_i}{n_{I_a}}k_a\quad(i\in I_a)
\end{equation}
give, on scalar coordinates with form
$\sum_i a_ib_i/n_i$, the orthogonal splitting
\[
 0\longrightarrow\bigoplus_aU_{I_a}
 \longrightarrow U_r\xrightarrow{G}U_s\longrightarrow0.
\]
Tensoring with $H$ gives the scalar flag-complex splitting.
Indeed $GJ=1$, and a within-group trace-zero vector pairs
trivially with $J(k)$. The induced pairing on the grouped
coordinates has denominator $n_{I_a}$. Successive sums compose,
and the section coefficients multiply as
$(n_i/n_{I_a})(n_{I_a}/n_J)=n_i/n_J$.
Thus the splittings agree under every refinement.

\subsection{A family with torsion in its flag cohomology}

Let $B=\mathbb C[t]$ and consider the rank-two Dolbeault operator
\[
 \bar\partial_t=\bar\partial+t\eta e_{12},\qquad
 e_{12}=\begin{pmatrix}0&1\\0&0\end{pmatrix}.
\]
It squares to zero because $\bar\partial\eta=\eta^2=0$.
It defines extensions
$0\to\mathcal O_X\to F_t\to\mathcal O_X\to0$.
They are algebraic in $t$: a \v{C}ech representative of the
elliptic extension class gives transition matrices
$1+t\eta_{ij}e_{12}$, whose products are finite.

The endomorphism model is
$M_2(H)\otimes B$ with differential $[t\eta e_{12},-]$.
For trace-free endomorphisms preserving the first summand, put
$e=e_{12}$ and $z=\operatorname{diag}(1,-1)$. Their complex is
\begin{equation}\label{cycompact:eq:nonsplit}
 \mathfrak b_t^0=(Hz\oplus He)\otimes B,
 \qquad dz=-2t\eta e,\quad de=0.
\end{equation}
The full endomorphism model retains associative composition;
the trace-free subcomplexes retain its Lie commutator.
To justify the model after twisting, filter matrices by the
diagonal weight for which $[e_{12},-]$ strictly raises weight.
The filtration is finite. On its associated graded, the
inclusion of \eqref{cycompact:eq:scalar-model} into Dolbeault
forms is a quasi-isomorphism. It follows that the twisted
inclusion is also a quasi-isomorphism. The same argument works
for upper-triangular endomorphisms and for any fixed finite
strictly upper-triangular matrix.

Directly from \eqref{cycompact:eq:nonsplit},
\begin{equation}\label{cycompact:eq:nonsplit-cohomology}
\begin{array}{c|c}
 j&H^j(\mathfrak b_t^0)\\\hline
 0&Be\\
 1&B\eta z\oplus(B/(t))\eta e\\
 2&B\sigma e\\
 3&B\eta\sigma z\oplus(B/(t))\eta\sigma e.
\end{array}
\end{equation}
The only nonzero differentials are multiplication by $-2t$
from $z$ to $\eta e$ and from $\sigma z$ to $\eta\sigma e$.
Every nonzero fibre has dimension one in each degree. The
derived zero fibre has dimension two in each degree, since
the differential of the free model becomes zero there.

Translations act trivially on $H^1(E,\mathcal O_E)$,
so for $t\ne0$ the extension on $E$ is invariant under
translation up to isomorphism. It has no translation
linearization. An $E$-equivariant vector bundle for the
free translation action on $E$ descends to a vector space
on a point and is therefore trivial. The nontrivial
extension has $h^0=1$: its connecting map
$H^0(\mathcal O_E)\to H^1(\mathcal O_E)$ is the nonzero
extension class. A trivial rank-two bundle has $h^0=2$.

More generally, let $V$ be a finite-dimensional matrix subspace
stable under $L=\operatorname{ad}N$, and put $\ell=\operatorname{rank}L$.
The complex $H\otimes V\otimes B$ with differential $t\eta L$
is the direct sum of the two-term complex $[V_B\xrightarrow{tL}V_B]$
and its shift by $[-2]$. Choosing bases for the image of $L$
gives
\[
 H^0=(\ker L)\otimes B,\qquad
 H^1\simeq(\operatorname{coker}L)\otimes B\oplus(B/(t))^\ell,
 \qquad H^{j+2}=\sigma H^j.
\]
The displayed splitting is a choice; its torsion submodule
and multiplicity are intrinsic. Derived specialization retains
$V$ in each of the four degrees.

\subsection{The endpoint traces in the flag pairing}

Let $N$ be strictly upper triangular, and use differential
$[\eta N,-]$ on
$\mathfrak g_r^0=\mathfrak{sl}_r\otimes H$ and on its
upper-triangular subcomplex $\mathfrak b_r^0$. Put
$V_r=H\otimes U_r$ for rank-one diagonal blocks, with zero
differential. Diagonal projection on upper-triangular matrices
is a chain map. The map
\[
 j:\mathfrak b_r^0\longrightarrow\mathfrak g_r^0\oplus V_r,
 \qquad x\longmapsto(x,\operatorname{diag}x)
\]
is isotropic for
\[
 Q((x,h),(y,k))=\int_X\operatorname{tr}(xy)-\int_X\sum_i h_ik_i.
\]
For upper-triangular $x,y$, the diagonal of $xy$ is the
componentwise product of their diagonals, proving isotropy.
The pairing on the target is perfect. In each coefficient
degree the matrix dimensions satisfy
$2(r(r+1)/2-1)=(r^2-1)+(r-1)$.
Pairing complementary upper and lower
matrix units then shows that the annihilator is exactly the
image of $j$. Hence
\begin{equation}\label{cycompact:eq:flag-pairing}
 0\longrightarrow\mathfrak b_r^0
 \xrightarrow{j}\mathfrak g_r^0\oplus V_r
 \longrightarrow(\mathfrak b_r^0)^\vee[-3]
 \longrightarrow0
\end{equation}
is an exact sequence of complexes, compatible with base change.
For $r=2$, omitting $V_r$ leaves
$\operatorname{tr}(z^2)=2$; the flag subspace would no longer
be isotropic. The endpoint trace term supplies the cancellation.

\section{Paired orbit directions and finite centres}
\label{cycompact:sec:orbit}

To preserve Calabi--Yau duality, removing an orbit tangent
direction also requires its paired obstruction direction.
The following complex
makes sense even when the orbit map changes rank. Let $R$ be
a commutative $\mathbb Q$-algebra, let $C$ be a finite-projective
complex in degrees zero through three, and equip it with a
perfect graded-symmetric pairing
$C\otimes_RC\to R[-3]$ compatible with its differential.
Let $W_R$ be finite projective and let $a:W_R\to C^1$ satisfy
$d_1a=0$. Define $a^\dagger:C^2\to W_R^*$ by
$(a^\dagger v)(w)=\langle aw,v\rangle$.

\begin{proposition}\label{cycompact:prop:orbit}
The complex
\begin{equation}\label{cycompact:eq:orbit}
 C_{\mathrm{orb}}=
 [C^0\oplus W_R\xrightarrow{(d_0,a)}C^1
 \xrightarrow{d_1}C^2
 \xrightarrow{(d_2,-a^\dagger)}C^3\oplus W_R^*]
\end{equation}
is perfect and self-dual with pairing of degree $-3$.
It commutes with derived base change, and
\[
 \det C_{\mathrm{orb}}=\det C\otimes(\det W_R)^2.
\]
A square root $P_C$ therefore gives the square root
$P_C\otimes\det W_R$.
\end{proposition}

\begin{proof}
Pairing compatibility and $d_1a=0$ imply $a^\dagger d_1=0$,
so consecutive differentials compose to zero. Add the evaluation
pairing between $W_R$ in degree zero and $W_R^*$ in degree
three. The two new terms in its differential are
$\langle aw,v\rangle-\langle w,a^\dagger v\rangle=0$.
The resulting pairing is perfect. Finite projectivity gives
derived base change and identifies both new determinant
contributions with $\det W_R$.
\end{proof}

Over a field, the outer terms record the kernel of the induced
orbit map $[a]:W_R\to H^1(C)$ and its dual. When $[a]$ is
injective, its image in $H^1(C)$ and its dual quotient in
$H^2(C)$ disappear. A cochain-level centre
$p:C^1\to W_R$ with $pa=1$ and $pd_0=0$ gives an actual
complement: its middle terms are $\ker p$ and $\ker a^\dagger$.
The projections $ap$ and $p^*a^\dagger$ split off two
contractible pairs in \eqref{cycompact:eq:orbit}. A splitting
only in the derived category gives the same conclusion by
splitting its homotopy idempotent in perfect complexes.

The rank-jumping example is already nonconstant over
$R=\mathbb C[t]$:
\begin{equation}\label{cycompact:eq:orbit-jump}
 [R\xrightarrow{t}R\xrightarrow{0}R\xrightarrow{-t}R].
\end{equation}
Its cohomology is $R/(t)$ in degrees one and three and zero
elsewhere. Its generic fibre is acyclic; its derived zero fibre
has one class in each degree zero through three. A fixed-rank
deletion cannot recover these fibres.

For $r$ disjoint point sheaves, the original Ext dimensions are
$(r,3r,3r,r)$. Removing the simultaneous orbit pair gives
$(r,3r-1,3r-1,r)$. Removing an orbit pair at every endpoint
would give $(r,2r,2r,r)$ and lose the relative contribution
\begin{equation}\label{cycompact:eq:orbit-excess}
 U[-1]\oplus U^*[-2],\qquad
 U=W^r/W_{\mathrm{diag}},\qquad
 \det\bigl(U[-1]\oplus U^*[-2]\bigr)=(\det U)^{-2}.
\end{equation}
All regroupings give coordinates on the same quotient $U$.
At the zero object, however, \eqref{cycompact:eq:orbit} is
$W[0]\oplus W^*[-3]$. After the tangent shift it is
$W[1]\oplus W^*[-2]$, the tangent of $T^*[-1]BE$ along its
zero section. A comparison with the relative unit $BE$ must
therefore specify what happens to its cotangent partner.

\subsection{A centre map retains a finite group}

\begin{proposition}\label{cycompact:prop:centre}
Let $M$ be an algebraic stack with an $E$-action and a map
$s:M\to E$ satisfying
$s(t\cdot m)=s(m)+Nt$ for an integer $N\ne0$.
Then, with $M_0=s^{-1}(0)$,
\begin{equation}\label{cycompact:eq:centre}
 [M/E]\simeq[M_0/E[N]].
\end{equation}
This equivalence preserves stabilizers.
\end{proposition}

\begin{proof}
The multiplication map $[N]:E\to E$ is finite \'{e}tale and
surjective. After that base change a point of $M$ admits
a translation into $M_0$. The choices differ exactly by
$E[N]$. Thus $E\times^{E[N]}M_0\to M$ is an isomorphism
after a faithfully flat base change and hence an isomorphism
of stacks. Quotienting by $E$ proves \eqref{cycompact:eq:centre},
with the same isomorphism groupoids on each side.
\end{proof}

For length-$N$ point sheaves, the sum of the elliptic support
cycle is such a map. On a bounded sector of fixed Hilbert
polynomial whose sheaves have finite image in $E$, choose an ample line bundle
$A$ on $S$ and a sufficiently large integer $m$.
The determinant of the perfect complex
$R\pi_{E,*}(F\otimes\pi_S^*A^m)$ is a line bundle on $E$.
Its divisor degree is the sum of the positive Hilbert
polynomials of the fibre-supported pieces. This gives a
positive integer $N$ and, after identifying $\operatorname{Pic}^N(E)$
with $E$, a centre map of the stated weight. The choice of
$m$ is made on that bounded sector.

Centred factors still have relative centres. For weights
$N_1,\ldots,N_r$, the condition for a centred total object
is $\sum_iN_ie_i=0$. Its solution is the full kernel of
$E^r\to E$, including finite components. Replacing this
kernel by a point would again erase extension parameters.
There is no positive-weight isogeny construction at $N=0$.

\section{A compact rational-curve coefficient}
\label{cycompact:sec:rational-curve}

Assume now that $S$ contains a smooth rational curve $C$.
Adjunction gives $C^2=-2$. A stable pair on $X$ consists of a
pure one-dimensional sheaf $F$ and a section
$\mathcal O_X\to F$ with zero-dimensional cokernel. Let
$P_n(X,([C],0))$ be its projective moduli scheme with
$\chi(F)=n$ and curve class $([C],0)$.

\begin{theorem}\label{cycompact:thm:rational-pairs}
For $r\geq0$, the universal effective divisor gives an
isomorphism
\[
 E\times\operatorname{Sym}^r C\xrightarrow{\sim}
 P_{r+1}(X,([C],0)),\qquad
 (e,D)\longmapsto
 (\mathcal O_X\longrightarrow i_{e,*}\mathcal O_C(D)).
\]
Here $i_e:C\hookrightarrow X$ is inclusion in the fibre over
$e$. Translation is free and its quotient is
$Y_r=\operatorname{Sym}^rC\simeq\mathbb P^r$.
The moduli scheme is empty for $n\leq0$.
\end{theorem}

\begin{proof}
An effective divisor in the class of $C$ must contain $C$,
since otherwise its intersection with $C$ would be nonnegative.
The remainder has zero class and hence zero intersection with
an ample divisor, so it vanishes. Thus $C$ is the unique
effective divisor in its class.

The degree against a positive divisor from $E$ is zero.
Every component of the sheaf's support cycle is consequently
constant in $E$. Projection to $S$ and the uniqueness just
proved show that the cycle is $C\times\{e\}$ with multiplicity
one. Its reduced ideal acts on $F$ with generically zero image.
That image is zero-dimensional, so purity makes it zero.
Hence $F=i_{e,*}L$ for a torsion-free rank-one sheaf $L$ on
the smooth curve $C$. It is a line bundle, and its nonzero
section gives $L=\mathcal O_C(D)$ for an effective divisor $D$.
Then $\chi(F)=\deg D+1$.

This classification gives the stated morphism in families.
To verify its scheme structure, use
$N_{C\times\{e\}/X}=\mathcal O_C(-2)\oplus\mathcal O_C$.
The local-to-global Ext sequence gives
\[
 0\to H^1(\mathcal O_C)\to\operatorname{Ext}^1_X(F,F)
 \to H^0(\mathcal O_C(-2)\oplus\mathcal O_C)
 \to H^2(\mathcal O_C),
\]
so $\operatorname{Ext}^1_X(F,F)=\mathbb C$.
Also $\operatorname{Hom}(F,F)=\mathbb C$,
$h^0(F)=r+1$ and $H^1(F)=0$. The tangent sequence for
deforming a sheaf with its section is
\[
 0\to\operatorname{Hom}(F,F)\to H^0(F)
 \to T_{(F,s)}P\to\operatorname{Ext}^1(F,F)\to H^1(F).
\]
The first map is injective because the section is nonzero.
Thus every tangent space has dimension $r+1$.

The morphism from $E\times\mathbb P^r$ is proper and
bijective on closed points, hence finite. It surjects onto
an irreducible target of dimension $r+1$. At each closed
point the local dimension equals the tangent dimension,
so the target is regular and normal. The finite map is
birational in characteristic zero and therefore an
isomorphism. Translation acts on the first factor.
\end{proof}

We can now specify an oriented critical enhancement on this
smooth quotient. Give $Y_r$ the zero function and take
$\mathcal Y_r=T^*[-1]Y_r$. Its virtual canonical line is
$K_{Y_r}^{\otimes2}$, with chosen square root
\begin{equation}\label{cycompact:eq:rational-orientation}
 L_r=K_{Y_r}=\mathcal O_{\mathbb P^r}(-r-1).
\end{equation}
This is a specified enhancement of the classical quotient.
An equivalence with a reduced derived stable-pair stack
requires a comparison of their derived structures.

Use the perverse vanishing-cycle normalization
$\phi_f^p=\phi_f[-1]$. For the zero function, the nearby-cycle
complex is zero. The defining triangle gives
$\phi_0F=F[1]$, and hence the oriented coefficient of
\eqref{cycompact:eq:rational-orientation} is
\begin{equation}\label{cycompact:eq:rational-coefficient}
 \varphi_r=\mathbb Q_{Y_r}[r],\qquad
 H_r=R\Gamma_c(Y_r,\varphi_r)
 \simeq\bigoplus_{k=0}^r\mathbb Qv_{r,k},
 \qquad |v_{r,k}|=2k-r.
\end{equation}
The differential is zero. The basis comes from powers of
the hyperplane class, since projective space has one cell
in each even real dimension. Exterior products of the
canonical lines give coherent orientations on ordered
products of these charts. Determinants of Jacobians
multiply under composition, so these product comparisons
agree under regrouping, with the usual Koszul signs.

There are two different functionals here. Integration is
\[
 \epsilon_r:H_r\longrightarrow\mathbb Q[-r],\qquad
 \epsilon_r(v_{r,r})=1,\quad
 \epsilon_r(v_{r,k})=0\quad(k<r).
\]
Its target has its nonzero term in degree $r$, the degree
of the top class. The categorical trace instead has domain
$\operatorname{End}^{\bullet}(H_r)$:
\[
 \tau_r(T)=\sum_j(-1)^j\operatorname{Tr}(T|_{H_r^j})
 \quad (|T|=0),\qquad \tau_r(T)=0\quad(|T|\ne0).
\]
For a formal degree-zero variable $y$, equip
$\prod_{r\geq0}\operatorname{End}^{\bullet}(H_r)$ with the
tail topology and componentwise composition. Define
$\mathcal T((T_r)_r)=\sum_{r\geq0}y^{r+1}\tau_r(T_r)$.

\begin{proposition}\label{cycompact:prop:rational-trace}
The maps $\tau_r$ and $\mathcal T$ are continuous cyclic chain maps, and
\begin{equation}\label{cycompact:eq:rational-trace}
 \tau_r(\operatorname{id})=(-1)^r(r+1),\qquad
 \mathcal T((\operatorname{id})_r)=\frac{y}{(1+y)^2}.
\end{equation}
The first number is also the Behrend-weighted Euler
characteristic of the quotient $Y_r$.
\end{proposition}

\begin{proof}
For homogeneous composable endomorphisms, pairing the two
matrix entries in opposite order gives
$\tau_r(AB)=(-1)^{|A||B|}\tau_r(BA)$.
In particular $\tau_r([d,T])=0$. This is the evaluation
trace of a perfect complex and is invariant under a change
of perfect model. All $r+1$ summands in
\eqref{cycompact:eq:rational-coefficient} have parity $r$.
Summing their signs and then the geometric-series derivative
gives \eqref{cycompact:eq:rational-trace}. Each coefficient
uses one finite complex, proving continuity. Finally the
Behrend function of the smooth scheme $\mathbb P^r$ is
$(-1)^r$, and its Euler characteristic is $r+1$.
\end{proof}

Oberdieck's quotient-incidence comparison identifies this
number with the reduced stable-pair incidence invariant for
the nonzero K3 class $[C]$ \cite{cycompact:Oberdieck}.
Upstairs, the smooth zero-function coefficient is
$\mathbb Q_E[1]\boxtimes\varphi_r$. Its restriction to the
slice over $0$, shifted by $[-1]$, is $\varphi_r$.
Its compact Euler characteristic vanishes because
$\chi(E)=0$, while the first quotient value is $+1$.
The shifted restriction and the quotient virtual-class
comparison specify how the nonzero invariant is obtained.

\section{Relative Ext and square roots along a curve}
\label{cycompact:sec:curve-ext}

The unframed sheaves
$G_e=i_*\mathcal O_C\boxtimes\mathcal O_e$, with
$i:C\hookrightarrow S$, test extension coefficients.
On $S$, the divisor resolution gives
\[
 R\operatorname{Hom}_S(i_*\mathcal O_C,i_*\mathcal O_C)
 \simeq\mathbb C\oplus\mathbb C[-2]
\]
as a complex of vector spaces: its two contributions are
$H^0(\mathcal O_C)$ and $H^1(\mathcal O_C(-2))$, the latter
shifted by one. The point resolution on $E$ then gives
Ext dimensions $(1,1,1,1)$ when the two parameters coincide.
For distinct parameters the supports are disjoint and
derived Hom vanishes.

For $m\geq1$, let $B_m=E^m/E_{\mathrm{diag}}\simeq E^{m-1}$, using the
first coordinate as origin. For $i\ne j$ let
$\delta_{ij}:B_m\to E$ be the difference $e_j-e_i$,
let $D_{ij}=\delta_{ij}^{-1}(0)$, and denote its canonical
divisor section by $s_{ij}$.

\begin{proposition}\label{cycompact:prop:curve-kernel}
After choosing the displayed splitting of the surface Hom
complex, the relative Hom from $G_{e_i}$ to $G_{e_j}$ is
represented by
\begin{equation}\label{cycompact:eq:curve-kernel}
 \mathcal K_{ij}=(\mathbb C\oplus\mathbb C[-2])\otimes
 [\mathcal O_{B_m}\xrightarrow{s_{ij}}
       \mathcal O_{B_m}(D_{ij})]_{0,1}.
\end{equation}
It computes derived restriction to every equality stratum,
including their nontransverse intersections, and all further
derived base changes.
\end{proposition}

\begin{proof}
On $E_z\times E_\delta$, resolve the fixed-point sheaf by
$[\mathcal O(-\{z=0\})\to\mathcal O]$.
Apply sheaf Hom into the graph sheaf $\mathcal O_{z=\delta}$.
Projection of that graph to $E_\delta$ is an isomorphism,
and the resulting complex is
$[\mathcal O_E\xrightarrow{s_0}\mathcal O_E(0)]_{0,1}$.
The section has a simple zero at $\delta=0$.
Translation and pullback by $\delta_{ij}$ give the elliptic
factor in \eqref{cycompact:eq:curve-kernel}. Tensoring with
the surface factor proves the formula. Its terms are vector
bundles, so it computes derived pullback along every map.
\end{proof}

On a formal normal disk, the elliptic factor is
$[\mathbb C[[t]]\xrightarrow{t}\mathbb C[[t]]]_{0,1}$.
It is quasi-isomorphic to its degree-one cokernel, but the
derived fibre of that cokernel has a degree-zero Tor term.
Both terms, and their shifts by $[-2]$, are needed for the
four collision groups. A coefficient pulled back from the
independent orbit quotient, which is a point, would have
constant derived fibres and cannot give this complex.

For a perfect complex use
$\det V^\bullet=\bigotimes_j(\det V^j)^{(-1)^j}$,
with parity its virtual rank modulo two. The two line
terms of \eqref{cycompact:eq:curve-kernel} occur in degrees
one and three; consequently
$\det\mathcal K_{ij}=\mathcal O(-2D_{ij})$ and its parity
is even. Choose constant determinant trivializations of the
diagonal self-Ext factors using the volume forms.
Taking the endomorphism determinant relatively over $B_m$ gives
\begin{equation}\label{cycompact:eq:curve-root}
 \det R\operatorname{Hom}_X
  \left(\bigoplus_iG_{e_i},\bigoplus_iG_{e_i}\right)
 =\mathcal O_{B_m}\left(-4\sum_{i<j}D_{ij}\right),\qquad
 \mathcal L_m=\mathcal O_{B_m}\left(-2\sum_{i<j}D_{ij}\right)
\end{equation}
is a specified square root.

For ordered blocks $A,B$, multiplication of divisor lines
gives
\begin{equation}\label{cycompact:eq:curve-root-sewing}
 \mathcal L_{A\sqcup B}\simeq
 r_A^*\mathcal L_A\otimes r_B^*\mathcal L_B\otimes
 \det\!\left(\bigoplus_{i\in A,\ j\in B}\mathcal K_{ij}\right),
\end{equation}
where $r_A,r_B$ retain the internal differences.
Every unordered pair of leaves is separated at exactly one
vertex of a binary tree, its least common ancestor. Each
parenthesization therefore counts the same divisor with
exponent $-2$. Multiplication of line bundles is associative,
and all cross determinants are even. Thus the maps themselves
agree under every ordered regrouping. They commute with
derived restriction, and $\mathcal L_0=\mathcal O$ supplies
the empty-block identity.

These roots are defined on labelled split parameter spaces.
At coincident parameters the direct sum acquires more
automorphisms. Descent under those automorphisms, extension
to nonsplit sheaves and comparison with a reduced critical
structure are further requirements. The divisor formula
does not specify a zero-function coefficient on $B_m$.

\section{Determinant sewing and orientation coherence}
\label{cycompact:sec:orientations}

For point sheaves there is a complementary global calculation.
Over $X^2$, let $\mathcal C$ be the relative derived Hom
between the two universal graph sheaves. It is perfect and
acyclic off the diagonal. Its determinant has a canonical
acyclic trivialization there. Since $X^2$ is smooth and its
diagonal has codimension three, this trivialization extends
uniquely: a rational section of a line bundle on a normal
variety with no zeros or poles in codimension one is regular
and invertible. Uniqueness makes it invariant under common
translation. The coefficient-line automorphisms act on its
determinant with exponent the Euler characteristic, which is
zero. It therefore descends to the common-torsor pair stack.

On the diagonal, the complex has cohomology
$\Lambda^*T_xX$ and dimensions $(1,3,3,1)$.
Its determinant is trivial because the exponent of
$\det T_xX$ is $-1+2-1=0$. The full complex still has
all four groups. For $r$ coincident copies its Yoneda algebra
is $M_r(\mathbb C)\otimes\Lambda^*T_xX$.
For independent tangent vectors $v_1,v_2$ and $r=2$,
the first-order deformation $E_{12}v_1+E_{21}v_2$ has
quadratic obstruction
$\operatorname{diag}(1,-1)v_1v_2\ne0$ although its scalar
trace is zero.

Let $0=F_0\subset F_1\subset F_2\subset F_3=F$ be a
family over a locally noetherian complex scheme $B$, with
flat length-one quotients $A_i$.
For projection $\pi:X\times B\to B$, put
\[
 C_{ij}=R\pi_*R\mathcal Hom(A_i,A_j),\qquad
 L_{ij}=\det C_{ij},\qquad
 D(F)=\det R\pi_*R\mathcal Hom(F,F).
\]

\begin{proposition}\label{cycompact:prop:flag-determinant}
There is a natural isomorphism
\begin{equation}\label{cycompact:eq:flag-determinant}
 D(F)\simeq\bigotimes_{i=1}^3D(A_i)
                \otimes\bigotimes_{i<j}L_{ij}^{\otimes2}.
\end{equation}
It is compatible with derived base change and common-torsor
descent. Both parenthesizations give the same determinant
map at nonsplit flags and collisions.
\end{proposition}

\begin{proof}
The graph sheaves are perfect, and the extensions make
the $F_i$ perfect. Locally choose bounded locally free
resolutions of the three quotients. Lifting their successive
extensions gives an upper-triangular differential on their
graded direct sum. Filtering the endomorphism complex by
matrix-block weight has associated graded the nine Hom
complexes $C_{ij}$. This statement uses the filtered
resolution, so it does not split the original extensions.

For a short exact sequence, the determinant isomorphism
concatenates a basis of the subobject with lifts of a basis
of the quotient. Changing the lifts leaves the top wedge
unchanged. Concatenating three quotient bases is independent
of parentheses. Degreewise application, using determinant
duals in odd degrees, gives the same refinement rule for
perfect complexes. Acyclic factors cancel canonically,
so filtered resolution comparisons preserve it.

Apply this rule to the nine blocks in the order
$C_{31},C_{21},C_{32},C_{11},C_{22},C_{33},C_{12},C_{23},C_{13}$.
Serre duality gives $C_{ji}\simeq C_{ij}^{\vee}[-3]$ and
therefore $\det C_{ji}=L_{ij}$. Every block has virtual
rank zero. The reorderings are thus even and give
\eqref{cycompact:eq:flag-determinant}. Both parenthesizations
refine to this same ordered filtration. Perfect base change
and the translation-invariant Serre pairing give the
remaining compatibilities.
\end{proof}

An orientation must lift this determinant identity to chosen
square roots. A line $D$ has a square root precisely when
its \'{e}tale Kummer class $\delta(D)\in H^2(-,\mu_2)$ vanishes.
After square roots have been chosen, two binary sewing
isomorphisms may differ by a sign. On a three-step flag,
the ratio of the two parenthesized sewing maps is a
$\mu_2$-valued associator. Its pentagon is a cocycle equation
on the full flag and overlap diagram. One must trivialize
that cocycle coherently. This is the extra associativity
condition in strong orientation data from
Definition~\ref{cyhall:def:strong-orientation}.

Ordinary orientation data on compact Calabi--Yau threefold
moduli are available \cite{JoyceUpmeier}.
Their existence does not identify the square roots after
orbit reduction or supply the required reduced flag maps.
Even a trivial determinant permits an incoherent choice:
multiply only the binary map for masses $(2,1)$ by $-1$,
and leave the others unchanged. All squares remain the
same, but the two maps on masses $(1,1,1)$ differ by $-1$.

\subsection{Two sign defects that can be solved explicitly}

Suppose the associator depends only on masses. Write it as
$a:\mathbb N^3\to\mu_2$, normalized to one if any argument
is zero, and assume its pentagon
\[
 a(n,r,s)a(m,n+r,s)a(m,n,r)
 =a(m+n,r,s)a(m,n,r+s).
\]
Define $\varepsilon(m,0)=\varepsilon(m,1)=1$ and
\begin{equation}\label{cycompact:eq:mass-correction}
 \varepsilon(m,n)=\prod_{j=1}^{n-1}a(m,j,1)\quad(n\geq2).
\end{equation}
Multiplying binary sewing by $\varepsilon$ changes the
associator to
\[
 a'(m,n,r)=a(m,n,r)
 \frac{\varepsilon(m,n)\varepsilon(m+n,r)}
      {\varepsilon(n,r)\varepsilon(m,n+r)}.
\]
The recursion gives $a'(m,n,1)=1$. Its pentagon with last
argument one gives $a'(m,n,r+1)=a'(m,n,r)$.
Induction from $r=0$ proves $a'=1$. The correction preserves
units and squares.

There is also an explicit finite-group calculation. If $G$
has odd order and $\alpha:G^2\to\mu_2$ is a normalized
two-cocycle with trivial coefficient action, put
$\beta(g)=\prod_{k\in G}\alpha(g,k)$.
Multiplying the cocycle identity over $k$ gives
\[
 \alpha(g,h)^{|G|}\beta(gh)=\beta(g)\beta(h).
\]
Odd order makes $\alpha=\delta\beta$.
For $G=(\mathbb Z/2)^2$, however,
$\alpha(g,h)=(-1)^{g_1h_2}$ is a cocycle and is not a
coboundary: coboundaries on an abelian group are symmetric,
and this one is not. Thus $E[N]$ has no constant sign
obstruction of this kind for odd $N$, while even torsion
can retain one. Neither calculation trivializes a gerbe
that varies over the sheaf moduli.

\section{A polystable spherical collection}
\label{cycompact:sec:polystable}

An Euler matrix alone does not produce a polystable object.
We construct three spherical sheaves with one polarization
and one Hilbert polynomial, so that every multiplicity vector
belongs to an actual polystable locus.

\begin{theorem}\label{cycompact:thm:triple}
There is a projective K3 surface with an ample class $H$,
$H^2=8$, and line bundles $L_i=\mathcal O_S(D_i)$,
$1\leq i\leq3$, such that
\[
 D_i^2=-8,\qquad D_iD_j=-4\ (i\ne j),\qquad HD_i=0.
\]
They are stable with common Hilbert polynomial $4m^2-2$,
and their Ext dimensions are
\[
 \dim\operatorname{Ext}^{\bullet}_S(L_i,L_j)=
 \begin{cases}(1,0,1),&i=j,\\(0,2,0),&i\ne j.
 \end{cases}
\]
Thus each $L_i$ is spherical, their Euler matrix is
\begin{equation}\label{cycompact:eq:Cartan}
 G=\begin{pmatrix}2&-2&-2\\-2&2&-2\\-2&-2&2\end{pmatrix},
\end{equation}
and every nonzero $F_{\mathbf n}=\bigoplus_iL_i\otimes\mathbb C^{n_i}$,
with $n_i\geq0$, is polystable.
\end{theorem}

\begin{proof}
In three hyperbolic planes in the K3 lattice choose
$e_i^2=f_i^2=0$, $e_if_i=1$, with different planes orthogonal.
Set
\[
\begin{split}
 D_1&=e_1-4f_1,\qquad D_2=-4f_1+e_2-4f_2,\\
 D_3&=-4f_1-4f_2+e_3-4f_3,\qquad V=f_1+f_2+2f_3.
\end{split}
\]
Their lattice $N$ has Gram matrix
\[
 \begin{pmatrix}
 -8&-4&-4&1\\-4&-8&-4&1\\-4&-4&-8&2\\1&1&2&0
 \end{pmatrix},\qquad\det N=-128.
\]
The minor in coordinates $e_1,e_2,e_3,f_1$ is one, so the
embedding is primitive. The first three generators span a
negative-definite subspace and their Schur complement is
positive; the signature is $(1,3)$.

Put $H=D_3+4V=e_3+4f_3$. Then $H^2=8$, $HD_i=0$ and
$HV=2$. The integral orthogonal complement of $H$ in $N$
is spanned by the $D_i$, with
\[
 \left(\sum_i a_iD_i\right)^2
 =-4\left(\sum_i a_i^2+\left(\sum_i a_i\right)^2\right).
\]
The bracket is even, so this complement has no vector of
square $-2$. The period and lattice embedding theorems
realize $N$ as the N\'{e}ron--Severi lattice of a projective
K3 surface. Since $H$ lies on no root wall, the Weyl chamber
description permits a marking in which it is ample
\cite[Chapters 8 and 14]{cycompact:Huybrechts}.

For $i\ne j$, the class $D_j-D_i$ is nonzero, has square
$-8$ and is orthogonal to $H$. Neither sign is effective,
so its line bundle has $H^0=H^2=0$. Riemann--Roch gives
$\chi=2-8/2=-2$, proving the cross-Ext calculation.
Self-Ext is $H^*(S,\mathcal O_S)$. Finally
$\chi(L_i(mH))=2+(D_i+mH)^2/2=4m^2-2$.
Rank-one torsion-free sheaves are Gieseker stable: a proper
same-rank subsheaf has a nonzero torsion quotient and a
strictly smaller reduced Hilbert polynomial. The common
polynomial makes the direct sums polystable.
\end{proof}

The Mukai vectors are $v_i=(1,D_i,-3)$ and satisfy
$-\langle v_i,v_j\rangle=G_{ij}$.
They are independent, so this gives an additive isometry
from their charge lattice with the negative Mukai form to
the lattice with matrix \eqref{cycompact:eq:Cartan}.
Budur--Zhang formality applies to every $F_{\mathbf n}$
\cite{cycompact:BudurZhang}. We will also need compatibility
with the geometric trace and with the whole collection.

\subsection{One cyclic model for all multiplicities}

Let $\mathcal D$ be the three-object Dolbeault dg category
with morphisms
$\mathcal A^{0,*}(S,L_i^\vee\otimes L_j)$ and differential
$\bar\partial$. Composition is wedge product and bundle
composition. Its degree-$-2$ trace pairing is
\[
 \langle a,b\rangle=\int_S\omega_S\wedge\operatorname{tr}(ba).
\]
Put $A(i,j)=\operatorname{Ext}^*_S(L_i,L_j)$ and
$\ell=\bigoplus_i\mathbb C1_i$.
A cyclic comparison preserves the associated Calabi--Yau
class; on a cyclic minimal model that class is represented
by this perfect pairing.

\begin{proposition}\label{cycompact:prop:cyclic-formality}
There is a strictly unital cyclic minimal model $A$ of
$\mathcal D$ with only its Yoneda product $m_2$ nonzero.
One cyclic quasi-equivalence $F=(F_n):A\to\mathcal D$
can be chosen for the entire three-object collection.
\end{proposition}

\begin{proof}
The collection is simple-minded: its negative Ext groups
vanish and $A^0=\ell$ consists of the vertex identities.
The cyclic minimal-model and strict-unit theorem for such
a $2$-Calabi--Yau collection applies, since Serre
duality is perfect of degree $-2$
\cite[Lemmas 4.4--4.5 and Proposition 4.8]{cycompact:DavisonPurity}.
It supplies a cyclic model with strict units over the
semisimple algebra $\ell$ and the given geometric trace.

The higher products then vanish for a direct reason.
For $n\geq3$, an input in $A^0=\ell$ makes $m_n$ zero
by strict unitality. If all inputs have positive degree
and one has degree two, the output has degree at least
three, since $|m_n|=2-n$. Only inputs all in degree one
remain; their output is in $A^2$. Pair that output with
each vertex unit. Cyclicity moves the unit into an input
of $m_n$, giving zero. The pairing between $A^2$ and
$A^0$ is perfect, so the output is zero.
\end{proof}

This is a formal statement in the completed tensor algebra
over $\ell$. The formal correction of a cyclic form can
be seen weight by weight. In the noncommutative de Rham
complex, give every coordinate and differential weight
one. For a closed error $\omega_w$ of weight $w>2$ beyond
the constant two-form,
Cartan's identity gives
$\omega_w=d_{\mathrm{dR}}(w^{-1}\iota_{\mathcal E}\omega_w)$,
where $\mathcal E$ is the Euler vector field.
The constant nondegenerate form identifies that primitive
with a vector field whose formal flow removes the error.
Increasing weights converge in tensor length. Analytic
convergence is a separate condition.

For multiplicity spaces $V_i$, amplify the morphisms to
$A(i,j)\otimes\operatorname{Hom}(V_i,V_j)$ and define
the Taylor coefficients of $F^{\mathrm{mat}}$ by summing
over composable matrix paths:
\begin{equation}\label{cycompact:eq:matrix-path}
 (F_n^{\mathrm{mat}}(a_n,\ldots,a_1))_{i_ni_0}
 =\sum_{i_1,\ldots,i_{n-1}}
 F_n((a_n)_{i_ni_{n-1}},\ldots,(a_1)_{i_1i_0}).
\end{equation}
The formula includes multiplication of the matrix
coefficients. Use the suspended bar convention on $A[1]$:
the bar differential has degree one and $F$ has degree
zero. Its tensor products determine the Koszul signs.

A finite twisted flag is an ordered direct sum with a
strictly upper-triangular degree-one matrix $\delta$
satisfying its Maurer--Cartan equation, here $\delta^2=0$.
Twisting inserts copies of $\delta$ into the operations;
upper-triangular morphisms preserve the filtration.
The cyclic class is carried by the full matrix morphism
complexes; its restriction to flag-preserving morphisms
need not be nondegenerate.

\begin{theorem}\label{cycompact:thm:cyclic-flags}
The fixed functor $F$ extends to a cyclic quasi-equivalence
on finite twisted objects, retaining their flag-preserving
subcomplexes. For morphisms $a_i$ between flags with
twisting matrices $\delta_i$, it is given by
\begin{align}
 F_*(\delta)&=\sum_{n\geq1}
    F_n^{\mathrm{mat}}(\delta,\ldots,\delta),
       \label{cycompact:eq:twisted-object}\\
 F_k^{\mathrm{tw}}(a_k,\ldots,a_1)
 &=\sum_{j_0,\ldots,j_k\geq0}
 F_{k+j_0+\cdots+j_k}^{\mathrm{mat}}
 (\delta_k^{j_k},a_k,\ldots,a_1,\delta_0^{j_0}).
       \label{cycompact:eq:twisted-morphism}
\end{align}
A power denotes repeated inputs. These maps are equivariant
under changes of multiplicity bases and coherent under
every ordered regrouping. They also apply over nilpotent
Artinian bases and in their inverse-limit formal completion.
\end{theorem}

\begin{proof}
Each upper-triangular input advances along a finite flag.
The object formula terminates before its length, and the
morphism formula terminates between each pair of its
distinguished inputs. With Artinian parameters, combine
flag distance with powers of the maximal ideal. Every
formula is then finite at each precision.

Evaluate the bar identity $b_{\mathcal D}F=Fb_A$ on repeated
copies of $\delta$. It gives the Maurer--Cartan equation
for $F_*(\delta)$. Marking $a_1,\ldots,a_k$ in the same
identity gives all twisted functor equations. A finite
filtration by matrix distance has associated graded a
direct sum of the original quasi-isomorphisms $F_1$.
Thus the twisted functor is quasi-fully faithful.

For essential surjectivity, induct on the flag length.
The last extension class lifts through quasi-full-faithfulness;
a change by a boundary gives an isomorphic twisted object.
Taking its cone extends the lifted flag. Over a square-zero
Artinian extension with kernel $J$, a tentative lift has
curvature in degree-two cohomology of the twisted
endomorphism complex tensored with $J$. The existing
target lift makes its image zero; the quasi-isomorphism
already proved on these complexes removes it.
Surjectivity on degree-one cohomology matches the extension class and
a degree-zero gauge change matches representatives.
Induction proves the Artinian claim.

The amplified trace is the original cyclic trace followed
by matrix trace. Its terms are closed matrix paths, whose
rotation has exactly the original Koszul sign. This also
transports the cyclic class after twisting. Basis changes
conjugate the path sums. Regrouping partitions the same
paths into consecutive subpaths, leaving their sum
unchanged. All subdivisions therefore give the required
coherences. The compatible finite-precision formulas pass
to the formal inverse limit.
\end{proof}

For three blocks the first nonlinear term is already
visible:
$ (F_*\delta)_{13}=F_1(\delta_{13})+
F_2(\delta_{12},\delta_{23})$,
with the composition order of \eqref{cycompact:eq:matrix-path}.
Separate choices for the two binary extensions need not
retain that term.

\section{The cubic critical model and its gauge directions}
\label{cycompact:sec:critical}

Choose a quiver $Q$ with vertices $1,2,3$ and two arrows
$i\to j$ for each $i<j$. Write $t(a),h(a)$ for the tail
and head of an arrow. With $\dim V_i=n_i$, its doubled representation space
has maps
$X_a:V_{t(a)}\to V_{h(a)}$ and
$Y_a:V_{h(a)}\to V_{t(a)}$.
Put $G_{\mathbf n}=\prod_iGL(V_i)$ and
$\mathfrak g_{\mathbf n}=\bigoplus_i\operatorname{End}(V_i)$.
The degree-one Yoneda products, using Serre-dual arrow
bases, give the moment map
\begin{equation}\label{cycompact:eq:moment}
 \mu_i=\sum_{h(a)=i}X_aY_a-\sum_{t(a)=i}Y_aX_a.
\end{equation}
The equation $x^2=0$ in the surface Ext algebra is
$\mu=0$. Its pointed formal gauge group over an Artinian
base is $\exp(\mathfrak g_{\mathbf n}\otimes\mathfrak m)$,
where $\mathfrak m$ is the maximal ideal. The full affine
quotient by $G_{\mathbf n}$ is the quiver moduli model.

For the embedding $S\times\{e\}\hookrightarrow X$, a
normal parameter normalized by $\omega_E$ gives the exterior Koszul factor
$\Lambda(\theta)$, with $|\theta|=1$ and
$\operatorname{tr}(\theta)=1$. Tensoring the formal cyclic
collection with this factor gives a formal cyclic
threefold comparison. A degree-one element is $x+z\theta$,
with $x$ in the surface degree one and $z$ in degree zero.
Its square is $x^2+(xz-zx)\theta$.

On
$\mathcal V_{\mathbf n}=\operatorname{Rep}(\overline Q,\mathbf n)
\oplus\mathfrak g_{\mathbf n}$ these equations are the
critical equations of
\begin{equation}\label{cycompact:eq:potential}
 \mathcal W_{\mathbf n}(X,Y,Z)
 =\sum_i\operatorname{Tr}(Z_i\mu_i)
 =\sum_a\operatorname{Tr}
       \bigl(Y_a(Z_{h(a)}X_a-X_aZ_{t(a)})\bigr).
\end{equation}
Differentiation with respect to $Z,X,Y$ gives, respectively,
\[
 \mu_i=0,\qquad
 Y_aZ_{h(a)}=Z_{t(a)}Y_a,\qquad
 Z_{h(a)}X_a=X_aZ_{t(a)}.
\]
Thus the classical critical quotient parametrizes
preprojective representations with a commuting normal
operator.

The derived critical quotient also keeps the gauge duals.
For a smooth affine $V$ with a smooth affine group $G$
acting and an invariant function $f$, write
$\rho:\mathfrak g\to\Gamma(V,T_V)$ for its action.
In \'{e}tale coordinates $x_i$, write
$\rho(a)=\sum_i\rho_a^i\partial_{x_i}$ for a basis element
$a\in\mathfrak g$. Introduce odd variables $\xi_i$
of degree $-1$ and variables $b_a$ of degree $-2$, and set
\begin{equation}\label{cycompact:eq:critical-dga}
 K_f=\operatorname{Sym}_{\mathcal O(V)}
       (T_V[1]\oplus(\mathfrak g\otimes\mathcal O(V))[2]),\qquad
 d\xi_i=\partial_if,\quad
 db_a=-\sum_i\rho_a^i\xi_i,\quad dx_i=0.
\end{equation}
The algebra is graded commutative. The action on all
generators is the induced $G$-action.

\begin{proposition}\label{cycompact:prop:critical-gauge}
The stack $[\operatorname{Spec}K_f/G]$ is the derived
critical locus of $f$ on $[V/G]$. At a critical point,
its tangent complex shifted to endomorphism degrees is
\begin{equation}\label{cycompact:eq:critical-tangent}
 [\mathfrak g\xrightarrow{\rho}T_V
 \xrightarrow{\operatorname{Hess}f}T_V^*
 \xrightarrow{-\rho^*}\mathfrak g^*]
\end{equation}
in degrees zero through three.
\end{proposition}

\begin{proof}
The tangent of $[V/G]$ pulled back to $V$ is
$[\mathfrak g\to T_V]$ in degrees $-1,0$.
The Koszul model for the intersection of its zero section
and $df$ in the cotangent stack has generators
$T_V[1]$ and $\mathfrak g[2]$, with the differential
\eqref{cycompact:eq:critical-dga}
\cite[Section 4]{cycompact:BeniniSafronovSchenkel}.
It squares to zero because $d^2b_a=-\rho(a)f=0$.
Linearization and the quotient tangent give
\eqref{cycompact:eq:critical-tangent}.
Differentiating invariance at a critical point gives
$\operatorname{Hess}(f)\rho=0$ and its dual identity.
This tangent construction makes the coordinate models
compatible on overlaps.
\end{proof}

For $N=\sum_i n_i>0$, set
$c=N^{-1}\sum_i\operatorname{Tr}Z_i$ and
$Z_i^0=Z_i-cI$. Let $\mathcal V_{\mathbf n}^0$ be the
hyperplane $\sum_i\operatorname{Tr}Z_i=0$ and
$\mathcal W_{\mathbf n}^0$ the restricted potential.
Since $\sum_i\operatorname{Tr}\mu_i=0$,
the potential is independent of $c$ and
\begin{equation}\label{cycompact:eq:critical-centre}
 \operatorname{Crit}_{[\mathcal V_{\mathbf n}/G_{\mathbf n}]}
       (\mathcal W_{\mathbf n})
 \simeq T^*[-1]\mathbb A^1\times
 \operatorname{Crit}_{[\mathcal V_{\mathbf n}^0/G_{\mathbf n}]}
       (\mathcal W_{\mathbf n}^0).
\end{equation}
At the algebra level the first factor is
$\mathbb C[c]\otimes\Lambda(\xi_c)$ with $d\xi_c=0$.
The centred factor retains all degree-$-2$ gauge generators.
This removes the orbit pair, whereas quotienting only
the translation of $c$ leaves $\xi_c$.

The local orientation is the canonical line of the smooth
quotient,
\[
 K_{[\mathcal V/G]}=\det(\mathcal V^*)\otimes\det\mathfrak g.
\]
Every arrow space pairs with its dual reverse-arrow space,
and the loop space is the adjoint representation. Their
determinant characters cancel. An ordering of the paired
spaces gives a trivialization. Removing the centre changes
this line by $\det W$, in agreement with
Proposition~\ref{cycompact:prop:orbit}.

\subsection{Dimensional reduction and the point chart}

For a function $f(x,z)=\sum_{a=1}^d z_af_a(x)$ linear in
the fibre coordinates, dimensional reduction identifies
compactly supported vanishing-cycle cohomology with
cohomology of $(f_a=0)$, shifted by $[-2d]$ and
Tate-twisted by $(-d)$ \cite[Appendix A]{DavisonCriticalCoHA}.
With $\varphi_f=\phi_f[-1]$, the constructed potential gives
\begin{equation}\label{cycompact:eq:dimred}
 R\Gamma_c(\mathcal V_{\mathbf n},\varphi_{\mathcal W}\mathbb Q)
 \simeq R\Gamma_c(\mu^{-1}(0),\mathbb Q)
       [-2\dim\mathfrak g_{\mathbf n}](-\dim\mathfrak g_{\mathbf n}).
\end{equation}
The equivariant version uses finite-dimensional
approximations. Centring replaces the loop dimension by
$\dim\mathfrak g_{\mathbf n}-1$; the zero locus stays the
same because $\mu$ already has total trace zero.
Equation~\eqref{cycompact:eq:dimred} is a direct-image
comparison, with the shift from compact affine fibres.

For a point sheaf there is an analytic cyclic model as
well. The volume germ identifies its cyclic minimal
algebra with $\Lambda^*T_xX$, its exterior product and
its top-degree trace \cite{cycompact:VanGarderen}.
With multiplicity $n$, matrix amplification gives
\begin{equation}\label{cycompact:eq:point-potential}
 f_n(X,Y,Z)=\operatorname{Tr}(Z[X,Y]).
\end{equation}
For an oriented tangent basis, if $a=Xe_1+Ye_2+Ze_3$,
then $f_n=\operatorname{vol}_x\operatorname{Tr}(a^3)/3$.
This fixes the normalization. Its critical equations
are the three commutator equations. In the formal
neighbourhood of the origin, or in a sufficiently small
analytic neighbourhood, they describe the corresponding
finite-length local-coordinate modules.

Writing the dual odd matrices as $X^*,Y^*,Z^*$ and the
degree-$-2$ gauge matrix as $b$, the full differential is
\[
 dX^*=[Y,Z],\quad dY^*=[Z,X],\quad dZ^*=[X,Y],\qquad
 db=-[X,X^*]-[Y,Y^*]-[Z,Z^*].
\]
The Jacobi identity gives $d^2b=0$. A single point retains
endomorphism dimensions $(1,3,3,1)$; after the paired
centre removal they are $(1,2,2,1)$.
The ambient orientation line is
$\det(\operatorname{End}\mathbb C^n)^{-2}$, trivialized by
the evaluation determinant. The local model retains the
scalar stabilizer as well as its dual.

\section{The first two-colour collision}
\label{cycompact:sec:two-colour}

For two one-dimensional vertex spaces, remove their common
normal centre and retain the relative coordinate $s$.
With $r$ arrows in each dual direction the potential is
\[
 f=sQ,\qquad Q(x,y)=\sum_{i=1}^r x_iy_i,
 \qquad (s,x,y)\in\mathbb A^{2r+1}.
\]
The effective relative group $\mathbb G_m$ has weights
$1,-1,0$ on $x,y,s$.

\begin{proposition}\label{cycompact:prop:collision}
The critical ideal is
\begin{equation}\label{cycompact:eq:collision-ideal}
 J=(Q,sx_1,\ldots,sx_r,sy_1,\ldots,sy_r)
   =(x,y)\cap(s,Q).
\end{equation}
For $r=2$, the affine quotient of $\{s=Q=0\}$ is
$\{u^2+vw=0\}$. Its quotient on $x\ne0$ is
$T^*\mathbb P^1$.
\end{proposition}

\begin{proof}
The derivatives give the first ideal. With
$\mathfrak m=(x,y)$, multiplication by $s$ is injective
modulo $\mathfrak m$. Since $Q\in\mathfrak m$, an element
$sa+Qb$ belongs to $\mathfrak m$ exactly when
$a\in\mathfrak m$, proving the intersection formula.
For $r=2$, the invariant matrix $(x_iy_j)$ has rank at
most one and trace zero. Write it as
$\left(\begin{smallmatrix}u&v\\w&-u\end{smallmatrix}\right)$;
its determinant is $-u^2-vw$. These entries generate
the invariant ring. On $x\ne0$, the vector $x$ defines
a line and $y(x)=0$ defines a cotangent vector there.
The reciprocal rescaling is precisely the cotangent
bundle quotient.
\end{proof}

For $r\geq2$ the two components are the spectral axis
and the irreducible collision quadric. The relative
gauge generator remains in the derived model, with
\[
 db=-\sum_i(x_i\xi_{x_i}-y_i\xi_{y_i}),
 \qquad d^2b=-\sum_i(x_isy_i-y_isx_i)=0.
\]

Set $\Phi_f=\phi_f[-1]\mathbb Q[2r+1]$ on the smooth
atlas. Its stalks detect more than the ordinary critical
locus.

\begin{proposition}\label{cycompact:prop:milnor}
For $r\geq2$, the stalk at the origin has one copy of
$\mathbb Q$ in degrees $1-2r,-1,0$ and no other
cohomology. At a nonzero point of the spectral axis
only degree $-1$ survives. At a smooth nonzero point
of the collision quadric only degree $1-2r$ survives.
\end{proposition}

\begin{proof}
The map $(s,x,y)\mapsto(s,x,sy)$ identifies $f^{-1}(1)$
with $\mathbb C^*\times Q^{-1}(1)$. Projection of
$Q^{-1}(1)$ to $x\ne0$ is an affine bundle with fibre
$\mathbb C^{r-1}$, so it has the homotopy type of
$S^{2r-1}$. Homogeneity identifies a small local Milnor
fibre with a sufficiently large ball sublevel in this
global fibre. Those sublevels stabilize: the squared
norm is proper and semialgebraic on the smooth affine
fibre, its critical set has finitely many connected
components, and its value is constant on each.
Above these critical values its gradient gives
deformation retractions. The Milnor fibre at the
origin therefore has type $S^1\times S^{2r-1}$.

Its reduced cohomology lies in degrees $1,2r-1,2r$.
The normalization gives
$H^j(\Phi_f)_0=\widetilde H^{j+2r}(F_0,\mathbb Q)$,
proving the first assertion. Away from the origin
on the axis the transverse function is the
nondegenerate quadratic form $Q$. Away from the
origin on the quadric use $Q$ as a coordinate;
the transverse function is the node $sQ$.
Their Milnor cohomology degrees are $2r-1$ and one,
respectively, giving the remaining shifts.
\end{proof}

For $r=1$, the two degree-one classes of the Milnor
fibre coincide in degree, giving a rank-two stalk
there. These are atlas coefficients. Passing to
the quotient requires their equivariant structures
and retains the classifying-stack directions.
Inverting $s$ cannot decide a class supported at
the collision: the complex
$[\mathbb C[s]e\xrightarrow{d}\mathbb C[s]]$ with
$|e|=-1$ and $de=s$ becomes contractible there,
but before localization has $H^0=\mathbb C[s]/(s)$.

\subsection{Real-root convolution relations}

In the constructible Hall algebra of nilpotent
representations of the preprojective algebra of $Q$,
use Euler-characteristic convolution. Let $e_i$ be
the characteristic function of the vertex simple.
The divided power $e_i^{(b)}=e_i^{*b}/b!$ is the
characteristic function of its $b$-fold direct sum,
since a complete flag in $\mathbb C^b$ has Euler
characteristic $b!$.

\begin{proposition}\label{cycompact:prop:Serre}
For distinct vertices $i,j$,
$(\operatorname{ad}e_i)^3(e_j)=0$.
Consequently the real positive Kac--Moody algebra
with matrix \eqref{cycompact:eq:Cartan} maps to
this constructible Hall algebra as a Lie algebra.
\end{proposition}

\begin{proof}
Evaluate the relation on dimension vector $(3,1)$
at $i,j$. Combine the forward arrows into
$X:\mathbb C^3\to\mathbb C^2$ and the reverse arrows
into $Y:\mathbb C^2\to\mathbb C^3$.
The preprojective relation $YX=0$ implies
$\operatorname{rank}X+\operatorname{rank}Y\leq2$.
Put $I=\operatorname{im}Y$ and $K=\ker X$.
A flag counted by $e_i^{(b)}*e_j*e_i^{(3-b)}$
is a $b$-plane with $I\subset U\subset K$.
If $I\not\subset K$, all terms are empty.
Otherwise write $r_0=\dim I$, $s_0=\dim K$.
The flag Euler characteristic is
$\binom{s_0-r_0}{b-r_0}$, and
$s_0-r_0=3-\operatorname{rank}X-\operatorname{rank}Y\geq1$.
Thus
\[
 \sum_{b=0}^3(-1)^b\binom{s_0-r_0}{b-r_0}
 =(-1)^{r_0}(1-1)^{s_0-r_0}=0.
\]
Multiplying by $3!$ gives the cubic Serre relation,
up to a common sign. Interchanging the vertices
proves the other relation.
\end{proof}

This constructs an actual homomorphism and its
enveloping-algebra map. It does not determine its
kernel or every imaginary primitive. Those require
the root-graded operators and their relations,
beyond the Euler matrix and the real-root Serre
identities.

\section{An oriented normal coefficient map}
\label{cycompact:sec:normal}

A local comparison can be made on the actual normal
Gysin map. Let $U$ be a smooth complex variety, let $m$
be a regular function, and put
\[
 p:Y=U\times\mathbb A^1_{\tau_{\mathrm g}}\to U,
 \qquad f=m\tau_{\mathrm g}.
\]
Let $j:U\hookrightarrow Y$ be the zero section,
$i:Z=(m=0)\hookrightarrow U$, and
$k:Z\times\mathbb A^1\hookrightarrow Y$.
The geometric coordinate $\tau_{\mathrm g}$ has degree
zero. Write $\Phi_f=\phi_f[-1]$, extended by zero to
$Y$; in particular $\Phi_0$ is the identity.

For $F\in D_c^b(U;\mathbb Q)$, normalize the complex
Thom map $j_*F\to p^*F[2]$ to have compact trace $+1$.
Applying vanishing cycles gives
\begin{equation}\label{cycompact:eq:normal-Gysin}
 G_F:j_*F\longrightarrow\Phi_fp^*F[2].
\end{equation}
The source identification follows from proper base change
and $fj=0$. Rank-one dimensional reduction, constructed
from the coefficient unit $p^*F\to k_*k^*p^*F$, gives
\[
 \theta_F:Rp_!\Phi_fp^*F[2]\xrightarrow{\sim}i_*i^*F.
\]
It applies to these pulled-back coefficients without a
smoothness assumption on $Z$.

\begin{lemma}\label{cycompact:lem:normal-trace}
With this orientation,
\begin{equation}\label{cycompact:eq:normal-trace}
 \theta_F Rp_!G_F=(F\longrightarrow i_*i^*F).
\end{equation}
The equality is natural in $F$ and in the normal frame.
\end{lemma}

\begin{proof}
Pull the Thom map to $Z\times\mathbb A^1$.
The Cartesian zero-section square identifies its
composite with restriction to $Z$ followed by the
same normal Thom map over $Z$. Its compact trace is
the identity. Naturality of the coefficient unit
before and after $\Phi_f$ proves
\eqref{cycompact:eq:normal-trace}. A complex frame
change preserves the real orientation and its
trace normalization.
\end{proof}

We now specify the coefficient change. Pull the
algebraic structure sheaf to the analytic site,
and write it as $\mathscr O_U$. Define the dg algebra
\begin{equation}\label{cycompact:eq:normal-Koszul}
 K_m=\mathscr O_U\otimes\Lambda(\eta,\chi),\qquad
 |\eta|=|\chi|=-1,\quad d\eta=m,\quad d\chi=0.
\end{equation}
It is the Koszul algebra of the two specified
equations $(m,0)$. The covariant exact realization is
\begin{equation}\label{cycompact:eq:normal-realization}
 \mathcal R_m(A)=K_m\otimes_{\mathbb Q_U}^{\mathbb L}Rp_!A,
 \qquad D_c^b(Y;\mathbb Q)\longrightarrow D(K_m).
\end{equation}
The target consists of sheaves of left dg modules;
its objects need not be coherent. This construction
includes extension of scalars by the specified
Koszul algebra. In unrestricted derived categories
the right adjoint forgets the module structure and
then applies $p^!$.

The algebraic normal resolution uses a separate
variable $\tau_{\mathrm a}$ of degree two. Set
$\mathcal E=\mathscr O_U[\tau_{\mathrm a}]\otimes\Lambda(e)$,
with $|e|=1$. Let $\lambda$ multiply by $e$ and
let $\eta$ contract with $e$. Thus
$\lambda\eta+\eta\lambda=1$, and
\[
 D_{\mathcal E}=m\lambda+\tau_{\mathrm a}\eta,
 \qquad D_{\mathcal E}^2=m\tau_{\mathrm a}.
\]
Put
\[
 \mathscr A=\operatorname{End}_{\mathscr O_U[\tau_{\mathrm a}]}
       (\mathcal E)\otimes\Lambda(\chi),\qquad
 \mathscr S=(\mathscr A[\tau_{\mathrm a}^{-1}]/\mathscr A)[-1].
\]
The differential of $\mathscr A$ is $[D_{\mathcal E},-]$;
on the shifted quotient it has the opposite sign.
Only finite negative Laurent sums are admitted.
Order $\eta,\chi$ before $\lambda$, and use the shifted
left action $a\cdot[x]=(-1)^{|a|}[ax]$.

For homogeneous $\xi\in K_m$ and $n\geq1$, define
\begin{equation}\label{cycompact:eq:normal-residue}
\begin{aligned}
 r(\xi\lambda\tau_{\mathrm a}^{-n})
 &=\begin{cases}(-1)^{|\xi|}\xi,&n=1,\\0,&n>1,\end{cases}\\
 r(\xi\tau_{\mathrm a}^{-n})&=0,\qquad
 s(\xi)=(-1)^{|\xi|}\xi\lambda\tau_{\mathrm a}^{-1}.
\end{aligned}
\end{equation}
They are degree-zero $K_m$-linear chain maps with
$rs=1$. The degree-$-1$ homotopy
\[
 h(\xi\tau_{\mathrm a}^{-n})
 =(-1)^{|\xi|+1}\xi\lambda\tau_{\mathrm a}^{-n-1},
 \qquad h(\xi\lambda\tau_{\mathrm a}^{-n})=0
\]
satisfies $dh+hd=1-sr$ and $h^2=rh=hs=0$.
Indeed $[D_{\mathcal E},\eta]=m$ and
$[D_{\mathcal E},\lambda]=\tau_{\mathrm a}$.
Substitution cancels the successive pole terms;
at pole one the regular term is discarded and
the remainder is exactly $1-sr$.

\begin{proposition}\label{cycompact:prop:normal-square}
The coefficient restriction
$q_F:K_m\otimes F\to K_m\otimes i_*i^*F$ is a
quasi-isomorphism. If
$\alpha_F=s\otimes1_F$ and
$\gamma_F=q_F^{-1}(1\otimes\theta_F)$ in the derived
category, then
\begin{equation}\label{cycompact:eq:normal-square}
 (r\otimes1_F)\alpha_F=\gamma_F\mathcal R_m(G_F).
\end{equation}
For constant coefficients, both realized objects
have the model $K_m$ and the realized arrow is the
identity. The oriented normal class is represented
by $\lambda/\tau_{\mathrm a}$, with residue $1$.
\end{proposition}

\begin{proof}
On $Z$, the coefficient restriction is the identity
on stalks. Off $Z$, multiplication by $\eta/m$
contracts the Koszul complex. Tensoring over the
field $\mathbb Q$ preserves these assertions,
so $q_F$ is a quasi-isomorphism.
Lemma~\ref{cycompact:lem:normal-trace} identifies
the composite of $\mathcal R_m(G_F)$ with
$1\otimes\theta_F$ as $q_F$.
Now $rs=1$ proves \eqref{cycompact:eq:normal-square}.
The resulting arrow is objectwise equivalent to
the identity of $K_m\otimes F$, giving the stated
models and zero cochain homotopy for their square.
\end{proof}

The geometric coordinate has been integrated out.
The algebraic coordinate has degree two, with the
Thom shift $[2]$ and the resolution shift $[-1]$
making $\lambda/\tau_{\mathrm a}$ degree zero.
These are distinct coordinate algebras. External
products of independent normal factors are compatible
by compact K\"unneth, the product orientation and
the Koszul tensor rule. The construction also glues
for a normal line bundle with a section of its dual.

\subsection{The character after fixing the framed slice}

Consider the framed matrices
\[
 A=\begin{pmatrix}w+u&B_0\\0&w\end{pmatrix},\quad
 B=\begin{pmatrix}y+\nu&D_0\\0&y\end{pmatrix},\quad
 \iota=\binom{\iota_1}{\iota_2},\quad j_{\mathrm{fr}}=(1,0).
\]
Their moment matrix is
\[
 [A,B]+\iota j_{\mathrm{fr}}
 =\begin{pmatrix}\iota_1&uD_0-\nu B_0\\\iota_2&0\end{pmatrix}.
\]
Eliminate the two regular framing equations by their
Koszul contractions and, on $B_0\ne0$, use diagonal gauge to set
$B_0=1$. The remaining equations are
$m=uD_0-\nu$ and zero, with generators $\eta,\chi$.
For an auxiliary matrix $L$, the trace potential is
$L_{11}\iota_1+L_{21}(uD_0-\nu B_0)+L_{12}\iota_2$.
Thus $\tau_{\mathrm g}=L_{21}$ is the normal coordinate.

Let $(s,t)\in(\mathbb C^*)^2$ scale $A,B$ by $s,t$,
and let $h_1,h_2$ denote its additive characters.
The compensating gauge is
$g=\operatorname{diag}(s^{-1},1)$ with framing change
$h=s^{-1}$. The transformations
\[
 A'=g(sA)g^{-1},\quad B'=g(tB)g^{-1},\quad
 \iota'=st\,g\iota h^{-1},\quad
 j'_{\mathrm{fr}}=h j_{\mathrm{fr}}g^{-1}
\]
preserve both slice conditions. They give
\begin{equation}\label{cycompact:eq:normal-weights}
\begin{array}{c|ccccc}
 \text{variables}&w,u&D_0&\nu,\eta&\chi&
       \tau_{\mathrm g},\tau_{\mathrm a},\lambda\\\hline
 \text{character}&h_1&h_2-h_1&h_2&h_1+h_2&-h_2.
\end{array}
\end{equation}
In fact $D_0'=(t/s)D_0$, $u'=su$, $\nu'=t\nu$,
and the lower-left entry of
$g((st)^{-1}L)g^{-1}$ is $t^{-1}\tau_{\mathrm g}$.
The zero $(2,2)$ moment component retains character
$h_1+h_2$. Hence the residue frame is invariant.
Under any degree-zero frame unit $a$,
\[
 \tau'_{\mathrm a}=a\tau_{\mathrm a},\quad
 \lambda'=a\lambda,\quad m'=a^{-1}m,\quad
 \eta'=a^{-1}\eta,
\]
and the residue and section agree in the two frames.
Writing $h_i$ also for the equivariant first Chern classes,
the normal Euler class is $-h_2$; compact Thom
normalization requires no inversion of that class.

\subsection{Supported specialization keeps the incidence term}

Set $A_0=\mathbb C[w,u]$, $z=w+u$ and $v=t_0-w$.
Use proper rational representatives at $t_0=\infty$
in the principal-parts modules
\[
\begin{split}
 P_z&=A_0[t_0,(t_0-z)^{-1}]/A_0[t_0],\\
 P_w&=A_0[t_0,(t_0-w)^{-1}]/A_0[t_0],\\
 P_{zw}&=A_0[t_0,((t_0-z)(t_0-w))^{-1}]/A_0[t_0].
\end{split}
\]
Let $\mathcal C=[P_z\oplus P_w\to P_{zw}]$ in degrees
$-1,0$, with differential $(a,b)\mapsto a-b$.
The target of the Cauchy comparison is
\[
 \mathcal T=\bigoplus_{p,r\geq1}\mathbb C[w]u^{-p}v^{-r}.
\]
Multiplication discards terms once either pole exponent
becomes nonnegative. This is a direct sum of finite
principal parts, with no completion in pole order.

The map $\Theta:\mathcal C\to\mathcal T$ is zero in
degree $-1$; in degree zero take first the $v$-principal
part at zero and then the $u$-principal part. Explicitly,
\begin{equation}\label{cycompact:eq:Cauchy}
 \Theta\!\left(\frac1{v^a(v-u)^b}\right)
 =(-1)^b\sum_{j=0}^{a-1}\binom{b+j-1}{j}
       u^{-b-j}v^{j-a},\qquad a,b\geq1.
\end{equation}
It is a quasi-isomorphism. The single-pole modules
intersect trivially by partial fractions after
inverting $u$ and torsion-freeness before inversion.
At each total pole order, the binomial expansion in
\eqref{cycompact:eq:Cauchy} is triangular with unit
diagonal, proving surjectivity. If a proper rational
function has zero image, its $w$-principal part has
no $u$-pole and lies in $P_w$. Subtract it; the
remainder has only a $z$-pole and lies in $P_z$.
Monic division by $t_0-z$ keeps its coefficients
in $A_0$, proving the kernel assertion.

Take $m=uD_0-\nu$ and put
\[
 M_\nu=\mathcal T[D_0,\nu]\otimes\Lambda(\eta,\chi),
 \qquad d\eta=uD_0-\nu,\quad d\chi=0.
\]
The residue contraction tensored with $\mathcal C$
and the Cauchy map give quasi-isomorphisms
\begin{equation}\label{cycompact:eq:supported-normal}
 \mathscr S\otimes_{A_0}\mathcal C
 \xrightarrow{r\otimes1}K_m\otimes_{A_0}\mathcal C
 \xrightarrow{1\otimes\Theta}M_\nu.
\end{equation}
They commute with derived specialization at $\nu=0$:
the Koszul terms are bounded free modules and the
principal-parts terms are free over $\mathbb C[\nu]$.
On the constructible side the vanishing-cycle
specialization need not be an isomorphism before
normal integration. After $Rp_!$, dimensional
reduction identifies it with the restriction of
the coefficient unit. The closed inclusion, that
unit and the relative compact trace commute with
base change, giving the same specialization after
extension to $K_m$.

Set
\[
 F_1=\frac1{(z-t_0)(t_0-w)},\quad
 c_1=((t_0-z)^{-1},(t_0-w)^{-1}).
\]
Then $dc_1=-uF_1$ and $\Theta(F_1)=u^{-1}v^{-1}$.
Before specialization the identities are
\begin{equation}\label{cycompact:eq:incidence-descendants}
\begin{split}
 d(\eta F_1+D_0c_1)&=-\nu F_1,\\
 d(\eta\chi F_1-\chi D_0c_1)&=-\nu\chi F_1.
\end{split}
\end{equation}
At $\nu=0$, the signed section lifts them to the
cycles
\[
\begin{split}
 B_\eta&=-\eta\lambda\tau_{\mathrm a}^{-1}F_1
                 +D_0\lambda\tau_{\mathrm a}^{-1}c_1,\\
 B_{\eta\chi}&=\eta\chi\lambda\tau_{\mathrm a}^{-1}F_1
                 +\chi D_0\lambda\tau_{\mathrm a}^{-1}c_1.
\end{split}
\]
Together with $\lambda\tau_{\mathrm a}^{-1}F_1$,
their images are
\begin{equation}\label{cycompact:eq:three-descendants}
 [u^{-1}v^{-1}],\quad[\eta u^{-1}v^{-1}],\quad
 [\eta\chi u^{-1}v^{-1}]
\end{equation}
in degrees $0,-1,-2$, and all are nonzero. For the
first class, evaluate a prospective boundary at
$D_0=0$. A degree-$-1$ boundary has only a $\chi$
component, so cannot give the second class. There
is no incoming differential in degree $-2$.
Multiplication by $u$ kills the displayed principal
part, proving closedness. Before projection,
omitting $D_0c_1$ leaves
$d(\eta F_1)=uD_0/[v(u-v)]\ne0$.

The specialization can also be read as Tor. On
$V=\mathcal T[D_0]$ let $N=uD_0$. The two-term
complex $[V[\nu]\xrightarrow{N-\nu}V[\nu]]$
has degree-zero cohomology $V$, with $\nu$ acting
as $N$, and no degree-$-1$ cohomology. Injectivity
follows by taking the highest $\nu$ coefficient;
polynomial division gives the cokernel. Its
derived zero fibre therefore retains
\[
 \operatorname{Tor}^{\mathbb C[\nu]}_1(V,\mathbb C)
 =\ker N=\bigoplus_{r\geq1}
       \mathbb C[w,D_0]u^{-1}v^{-r}.
\]
The cokernel of $N$ is $\mathcal T$, since $u$
acts surjectively on $\mathcal T$. Adjoining
$\chi$ retains both copies with their shifts.
On a generic fibre, however,
$(N-\nu)^{-1}=-\sum_{j\geq0}N^j\nu^{-j-1}$;
the sum is finite on each principal part.
The generic complex is contractible, while its
derived special fibre has the classes
\eqref{cycompact:eq:three-descendants}.

\subsection{The kernel test for a Hall comparison}

The functor \eqref{cycompact:eq:normal-realization}
has a concrete nonzero kernel. Let
$a:U\times\mathbb C^*\hookrightarrow Y$ and put
$L=Ra_*\mathbb Q_{U\times\mathbb C^*}$.
On each normal line the localization triangle is
\[
 \mathbb Q_{\{0\}}[-2]\xrightarrow{\mathrm{Thom}}
 \mathbb Q_{\mathbb C}\longrightarrow
 Ra_*\mathbb Q_{\mathbb C^*}\longrightarrow.
\]
Compact pushforward takes its first arrow to the
identity of $\mathbb Q[-2]$. Hence $Rp_!L=0$,
although $L\ne0$. Its restriction to the zero
section is $\mathbb Q_U\oplus\mathbb Q_U[-1]$.
For the ordinary internal tensor,
\begin{equation}\label{cycompact:eq:normal-tensor-test}
 \mathcal R_m(L)=0,\qquad
 \mathcal R_m(L\otimes j_*\mathbb Q_U)
 =K_m\oplus K_m[-1].
\end{equation}
When $Z$ is nonempty the second object is nonzero.
Thus this realization is not strong monoidal for
that tensor. Its external-product comparison
remains valid.

There is a separate proper-image test on the
algebraic models. For $f:\mathbb P^1\to\mathrm{pt}$,
take $m=0$ and $K_0=\mathbb C\otimes\Lambda(\eta,\chi)$.
Constructible proper pushforward followed by
realization gives $K_0\oplus K_0[-2]$, since
$H^*(\mathbb P^1,\mathbb Q)$ has degrees zero
and two. Ordinary coherent pushforward of the
model $\mathcal O_{\mathbb P^1}\otimes K_0$
gives $K_0$. Indeed the degree-one \v{C}ech
quotient is
\[
 \mathbb C[z,z^{-1}]/
       (\mathbb C[z]+\mathbb C[z^{-1}])=0.
\]
The two complexes have cohomological widths four
and two; a uniform shift cannot identify them.
This excludes that proposed use of ordinary
coherent direct image. It is not a statement
about every functor on the sheaf-module target.

For an actual critical Hall convolution
$\star_{\mathrm H}$ and specified cut realizations
$\mathcal R_d$, descent through their kernels
$\mathcal N_d$ requires
\begin{equation}\label{cycompact:eq:Hall-kernel}
 \mathcal R_{d+e}(A\star_{\mathrm H}B)=0
 \quad\text{if }A\in\mathcal N_d\text{ or }B\in\mathcal N_e.
\end{equation}
The ordinary tensor test does not decide this
different product. The normal square constructs
its stated coefficient arrow; a Hall lift must
still prove \eqref{cycompact:eq:Hall-kernel} and
compare the actual proper traces and supported
incidence maps.

\section{Vertical, wrapped and mixed support}
\label{cycompact:sec:wrapped}

The projection $\pi_E:X\to E$ gives a useful
support decomposition. Call a coherent sheaf
\emph{vertical} if its support has finite image
in $E$, and let $\mathcal T$ be their full
subcategory. Let $\mathcal F$ consist of coherent
sheaves flat over $E$.

\begin{proposition}\label{cycompact:prop:torsion-pair}
The pair $(\mathcal T,\mathcal F)$ is a torsion
pair in $\operatorname{Coh}(X)$. Every sheaf has
a functorial exact sequence
\begin{equation}\label{cycompact:eq:torsion-pair}
 0\longrightarrow t_E(F)\longrightarrow F
 \longrightarrow F/t_E(F)\longrightarrow0
\end{equation}
with the first term vertical and the last flat
over $E$, and
$\operatorname{Hom}(\mathcal T,\mathcal F)=0$.
\end{proposition}

\begin{proof}
The sum of all vertical subsheaves stabilizes
because $F$ is noetherian; its support has finite
image in $E$. This defines $t_E(F)$.
The quotient has no torsion over the local
discrete valuation rings of the smooth curve
$E$, so it is flat over them. A map from a
vertical sheaf to an $E$-flat sheaf has
torsion image and is zero. These characterizations
also prove functoriality.
\end{proof}

Sheaves pulled back from $S$ are examples of
wrapped support, and belong to $\mathcal F$.
The torsion-pair sequence does not yield a
product decomposition of families. On $E$,
let $F_t$ be the family of extensions
\[
 0\to\mathcal O_E(-p)\to F_t\to\mathcal O_p\to0
\]
with extension parameter $t$.
For $t\ne0$, $F_t\simeq\mathcal O_E$ has no
torsion; at zero it is
$\mathcal O_E(-p)\oplus\mathcal O_p$ and its
torsion has length one. The total family is
flat in $t$, but its torsion parts are not.
Pushing it into $\{s\}\times E$ gives the same
example on $X$. Hall identities using
\eqref{cycompact:eq:torsion-pair} must therefore
use the corresponding strata of the extension
stack.

The mixed Ext groups retain the normal direction.
For $i:S\times\{e\}\hookrightarrow X$ and
perfect $A,B$ on $S$,
\begin{equation}\label{cycompact:eq:mixed-Ext}
\begin{split}
 R\operatorname{Hom}_X(\pi_S^*B,i_*A)
   &\simeq R\operatorname{Hom}_S(B,A),\\
 R\operatorname{Hom}_X(i_*A,\pi_S^*B)
   &\simeq R\operatorname{Hom}_S(A,B)\otimes W[-1],\\
 R\operatorname{Hom}_X(i_*A,i_*B)
   &\simeq R\operatorname{Hom}_S(A,B)
                  \otimes(\mathbb C\oplus W[-1]),\\
 R\operatorname{Hom}_X(\pi_S^*A,\pi_S^*B)
   &\simeq R\operatorname{Hom}_S(A,B)
                  \otimes R\Gamma(E,\mathcal O_E).
\end{split}
\end{equation}
The first two follow from the point resolution
on $E$: $R\operatorname{Hom}_E(\mathcal O_E,\mathcal O_e)
=\mathbb C$ and
$R\operatorname{Hom}_E(\mathcal O_e,\mathcal O_E)
=W[-1]$. The other two follow from its
self-Ext and from K\"unneth. Thus the two
directions of a mixed extension carry
different shifts, and both are needed in a
cyclic comparison.

\section{Reduced invariants and the remaining coefficient map}
\label{cycompact:sec:global-comparison}

A reduced virtual-class theory exists
for stable pairs on $S\times E$.
Let $\beta\ne0$ be an effective K3 curve class,
$d\geq0$, and put $P=P_n(X,(\beta,d))$.
Its standard perfect obstruction theory
$\mathbb E\to L_P$ has amplitude $[-1,0]$.
The semiregularity morphism is
\[
 \operatorname{sr}:\mathbb E^\vee
 \longrightarrow H^{1,3}(X)\otimes\mathcal O_P[-1].
\]
It is surjective on obstruction sheaves. Its
dual cone gives a reduced obstruction theory
of virtual dimension one. The existence of
the map from that cone to $L_P$ uses the
twistor deformation of the K3 surface; it
does not follow merely by deleting a vector
from the obstruction space.

The elliptic orbit vector and semiregularity
are Serre dual. One way to see their
relationship is to express both through the
universal Atiyah class: contraction with
the invariant vector field gives the orbit
direction, while composition and trace give
semiregularity. This duality allows a second
cone, removing the orbit direction, to carry
a symmetric obstruction theory on an actual
centre fibre. Determinant and Fourier--Mukai
constructions give a centre map of nonzero
weight, so that this fibre has a finite-group
quotient as in Proposition~\ref{cycompact:prop:centre}.

For the incidence convention, let $\mathcal I$ be the
universal stable-pair complex on $X\times P$, with
projections $\pi_X,\pi_P$. For $\gamma\in H^4(X,\mathbb Q)$
put
\[
 \tau_0(\gamma)=-\pi_{P,*}
       \bigl(\operatorname{ch}_2(\mathcal I)\cup\pi_X^*\gamma\bigr).
\]
This is a degree-two class on $P$.

\begin{theorem}[Reduced stable-pair comparison]
\label{cycompact:thm:reduced-pairs}
For these nonzero K3 classes, the quotient
$P/E$ has a zero-dimensional virtual class
whose pullback along $\pi:P\to P/E$ is
\begin{equation}\label{cycompact:eq:virtual-comparison}
 \pi^*[P/E]^{\mathrm{vir}}=[P]^{\mathrm{red}}.
\end{equation}
Its degree is the Behrend-weighted Euler
characteristic of the quotient, with finite
stabilizers counted orbifoldly. It equals
the reduced incidence invariant obtained by
inserting $\tau_0(\beta^\vee\boxtimes[0])$,
where $\langle\beta,\beta^\vee\rangle=1$.
\end{theorem}

This is the quotient and incidence theorem
of Oberdieck \cite[Theorems 1--2]{cycompact:Oberdieck}.
Concretely, if $K$ is a centre fibre with
finite group $G$, its class on the quotient
is $|G|^{-1}q_*[K]^{\mathrm{vir}}$.
The symmetric obstruction theory computes
its weighted Euler characteristic. The
virtual-class comparison uses the common
Fulton class and Segre-class expression;
the normalized incidence intersects the
elliptic orbit with degree one. This
construction is independent of the chosen
centre fibration.

The original standard virtual class is zero
in this sector. Likewise, a surjective
cosection has empty degeneracy locus for
the cosection-localized standard class
\cite[Theorem 1.1]{cycompact:KiemLi}.
The reduced class above uses the
separate reduced obstruction theory.
For $\beta\ne0$, the reduced Donaldson--Thomas
and stable-pair series also agree
\cite[Theorem 3(i)]{cycompact:Oberdieck}.
That comparison uses the local Hall identity
and Behrend integration for the quotient;
the vanishing of $\chi(S\times E)$ alone
does not construct it.

Equivariant Hall integration gives another
precise numerical construction. For the
elliptic translation action on the smooth
projective Calabi--Yau threefold, the
regular equivariant Hall algebra has a
semiclassical Poisson quotient. Its
integration map takes values in the
charge algebra with dual numbers
$\mathbb Q[\varepsilon]/(\varepsilon^2)$
\cite[Theorem 5]{OberdieckShen}.
The two terms integrate the fixed locus
and the moving quotient, respectively,
with the sign $(-1)^{\dim E}$ on the
latter. With $\chi$ the sheaf Euler pairing,
the target basis $c_\alpha$ has
product
$c_\alpha c_\beta=(-1)^{\chi(\alpha,\beta)}c_{\alpha+\beta}$
and Poisson bracket
$\{c_\alpha,c_\beta\}=(-1)^{\chi(\alpha,\beta)}
\chi(\alpha,\beta)c_{\alpha+\beta}$.

Write the integration as $I_0+\varepsilon I_1$.
Then $I_0$ is unital and multiplicative, and
\begin{equation}\label{cycompact:eq:dual-number}
 I_1(ab)=I_0(a)I_1(b)+I_1(a)I_0(b),\qquad
 I_1(1)=0.
\end{equation}
Thus the moving contribution is an
$I_0$-derivation. It supplies reduced
numbers, with a different domain and unit
law from a state-valued Hall action.

The preceding constructions fix substantial
parts of a global coefficient comparison:
the common torsor and its relative unit;
the disjoint coefficient and trace; the
full scalar and orbit complexes; explicit
compact coefficients and split roots;
one formal cyclic comparison on finite
flags; and an oriented normal Gysin map
with its supported specialization.
Their global assembly has a definite
remaining equation. For a proposed reduced
critical enhancement of the common-torsor
extension stack, write
$\operatorname{Filt}_{\alpha,\beta}$ for its
filtered objects, let $\operatorname{gr}$ take its
graded pair and let $\operatorname{ev}$
forget the filtration. For specified
virtual index $\iota_{\alpha,\beta}$,
the desired critical coefficient map is
\begin{equation}\label{cycompact:eq:global-Hall-map}
 \operatorname{gr}^*
    (\varphi_\alpha\boxtimes_{BE}\varphi_\beta)
       [\iota_{\alpha,\beta}]
 \longrightarrow
 \operatorname{ev}^!\varphi_{\alpha+\beta}.
\end{equation}

\begin{problem}\label{cycompact:prob:global}
Construct the compact reduced critical
coefficient system $\varphi$ with its
relative unit, cyclic pairing and strong
orientation, and the maps
\eqref{cycompact:eq:global-Hall-map} on the
full extension stack. Prove their
compatibility with every flag refinement,
finite-isogeny descent, collision
specialization and the support conditions
needed for the chosen pushforward.
Compare their trace and orbit reductions
with \eqref{cycompact:eq:flag-pairing} and
\eqref{cycompact:eq:orbit}, and their
normal restriction with
\eqref{cycompact:eq:supported-normal}.
\end{problem}

The ordinary critical Hall theorem of
Chapter~\ref{cyhall:sec:arrow} supplies its
coefficient map under the stated
Calabi--Yau, geometric and strong
orientation hypotheses. Applying it to
this reduced relative stack requires
those data on that stack. A further
comparison is then needed with the
independently defined chiral states:
it must preserve normal oscillators,
mixed wrapped collisions, the local
vacuum and cycle joining. The numerical
trace \eqref{cycompact:eq:rational-trace}
and integration \eqref{cycompact:eq:dual-number}
remain positive consequences in their
own coefficient categories.
